\section*{Chapter \ref{ch:ll:managed-and-functional-languages-on-the-desk} --- Managed and Functional Languages on the Desk}

\begin{solution}{exo:ll:managed-and-functional-languages-on-the-desk:1}
$64 \times 100\,000 = 6.4$~MB a second fills 1~GiB in 167.8 seconds: about 139 young collections in a 6.5-hour session
(\cref{prop:ll:managed-and-functional-languages-on-the-desk:budget}).
\end{solution}

\begin{solution}{exo:ll:managed-and-functional-languages-on-the-desk:2}
$0.001 \times 200\,000 = 200$ messages; the last arrives just before the pause ends and the first waits about a millisecond,
and all of them are then processed late in turn, so the queue takes a further 200 processing times to drain.
\end{solution}

\begin{solution}{exo:ll:managed-and-functional-languages-on-the-desk:3}
On the committed measurement the maximum falls from about \qty{19.6}{\milli\second} to \qty{0.58}{\milli\second}, some
thirty times, and the p99.9 from about \qty{26}{\micro\second} to \qty{1.7}{\micro\second}; the median is slightly worse,
\qty{335}{\nano\second} against \qty{313}{\nano\second}.
\end{solution}

\begin{solution}{exo:ll:managed-and-functional-languages-on-the-desk:4}
CPython triggers a young collection when allocations of tracked objects exceed deallocations by the threshold (700 by
default on Python 3.10). A bounded deque frees an old order for every new one, so the difference never grows; a growing
book adds more than it frees and crosses the threshold every few hundred messages.
\end{solution}

\begin{solution}{exo:ll:managed-and-functional-languages-on-the-desk:5}
The hot path is interpreted or lightly compiled until the just-in-time compiler has seen enough executions; the open is
also when rare paths are first taken and may deoptimise. Remedies: warm up by replaying recorded traffic through the real
code against a simulator before the open; keep types and branches on the hot path stable; or compile ahead of time.
\end{solution}

\begin{solution}{exo:ll:managed-and-functional-languages-on-the-desk:6}
$20 \times 10^6 / 6 \times 10^6 \approx 3.33$ seconds of input at full rate.
\end{solution}

\begin{solution}{exo:ll:managed-and-functional-languages-on-the-desk:7}
Freezing moves every object alive at that moment to a permanent generation that collections ignore, so full collections
scan only the orders added since the last freeze and are much shorter. The price is memory: frozen objects that die are
never reclaimed until the heap is unfrozen, so a book whose orders are cancelled keeps their memory.
\end{solution}

\begin{solution}{exo:ll:managed-and-functional-languages-on-the-desk:8}
Under a millisecond is a thousand microseconds, far beyond a microsecond budget; and a concurrent collector still takes
cores, memory bandwidth and cache from the program while it runs, and pauses the threads at safepoints. It affects latency
unless the gateway allocates nothing during the session.
\end{solution}

\begin{solution}{pb:ll:managed-and-functional-languages-on-the-desk:1}
\begin{enumerate}
\item $48 \times 200\,000 = 9.6$~MB a second; 224.64~GB in a session.
\item Every $2~\text{GiB}/9.6~\text{MB/s} = 223.7$ seconds: about 105 times a session.
\item $2~\text{GiB}/(200\,000 \times 23\,400) = 0.459$ bytes.
\item Zero: no allocation per message on the hot path.
\item $105 \times \qty{0.5}{\milli\second} \approx \qty{52}{\milli\second}$.
\item $0.0005 \times 200\,000 = 100$ messages; the first waits about \qty{0.5}{\milli\second}, and the queue then takes a
further 100 processing times to drain.
\item The allocation rate triples: about 314 collections and \qty{157}{\milli\second} paused, with 300 messages behind each
pause.
\item Queues and pauses scale with traffic, and the busiest days are the ones where latency is worth most.
\item Primitive fields and primitive collections (arrays of \texttt{long}), with no boxing.
\item A binary log: the hot path writes a format identifier and raw arguments to a pre-allocated ring; formatting happens
elsewhere, later.
\item Off the heap or in flat arrays of primitives allocated at start-up, so that no collection scans or moves them.
\item Measure allocated bytes per thread before and after a replay of messages (the JVM's per-thread allocation counters, or
an allocation profiler) and assert zero after warm-up.
\item \textbf{Named result.} The engine may allocate at most 0.459 bytes per message, in practice nothing, to collect once a
session; at 48 bytes it collects about 105 times and spends about \qty{52}{\milli\second} paused.
\item Safepoints and the collector's background threads, which still take cores and cache; and the start-up and warm-up
cost.
\item Because the just-in-time compiler still needs to see the hot path executed before it is optimised.
\item That removing per-message allocation removes the collector from the loop, and that the worst message is the full
collection over the live heap.
\item Readable, correct, efficient code, written and changed quickly: correctness per engineer more than speed per message.
\item When the budget is in the low microseconds at a high percentile, and the team cannot guarantee zero allocation in
every release.
\item The \omterm{def:ll:where-latency-comes-from:budget}{latency budget} at the percentile that matters, the allocation per message, and the measured pauses under a
replayed busy day.
\item A pause proportional to what the hot path allocates and to what the heap keeps alive.
\end{enumerate}
\end{solution}

\begin{solution}{iq:ll:managed-and-functional-languages-on-the-desk:1}
It traces reachable objects from the roots and reclaims the rest. Most objects die young, so new objects go to a small young
generation collected often, cheaply, in proportion to survivors; survivors are promoted to an old generation collected
rarely.
\iqlookfor{the weak generational hypothesis and cost in proportion to survivors.}
\end{solution}

\begin{solution}{iq:ll:managed-and-functional-languages-on-the-desk:2}
Allocate nothing per message after warm-up: pre-allocated pools and ring buffers, primitive arrays instead of boxed
collections, \omterm{def:ll:managed-and-functional-languages-on-the-desk:offheap}{off-heap memory} for large state, binary logging; a quiet collector and a young generation sized for the
residual; warm-up before the open; measure allocation in tests.
\iqlookfor{no allocation, off-heap state, warm-up, measurement.}
\end{solution}

\begin{solution}{iq:ll:managed-and-functional-languages-on-the-desk:3}
The time before the JIT has compiled the hot path with its profile; before the open, replay recorded traffic through the
real code (orders to a simulator) until the hot methods are compiled, and keep the profile representative.
\iqlookfor{warm-up with realistic traffic, never with real orders.}
\end{solution}

\begin{solution}{iq:ll:managed-and-functional-languages-on-the-desk:4}
The runtime discards optimised code whose assumptions (a single type at a call site, a branch never taken) have been broken,
and falls back to slower code; the first rare event of the day, often at a volatile moment, is what breaks them.
\iqlookfor{speculative optimisation and its fall-back.}
\end{solution}

\begin{solution}{iq:ll:managed-and-functional-languages-on-the-desk:5}
Correctness and speed of development: expressive types that encode invariants (every case of an order state handled), immutable
data by default, concise code that traders can review, as the firm that published its experience in 2008 argued; with
performance good enough for its budget.
\iqlookfor{correctness and review, backed by the published account.}
\end{solution}

\begin{solution}{iq:ll:managed-and-functional-languages-on-the-desk:6}
Record every collection's duration with the \texttt{gc} module's callbacks and correlate them with the slow messages; count
collections by generation. Then stop allocating per message (pre-allocated arrays), freeze long-lived objects after start-up,
or schedule full collections explicitly in quiet moments.
\iqlookfor{measure the collector, then remove its cause.}
\end{solution}
